% TOP_PROBLEM: {"id": 30007099, "title": "Periodic translational tiling in dimension three", "category_id": 2, "category": {"id": 2, "name": "combinatorics", "display_name": "Combinatorics", "description": "Counting problems, graph theory, discrete structures.", "slug": "combinatorics", "order_index": 2, "created_at": "2026-07-31T15:26:25.671Z"}, "status": "open"}
% FIRST_POSTED: 2026-10-06
% ATTEMPT_STATUS: solved
% OPENAI_CLAIM: {"schema_version":1,"source_repo":"https://github.com/openai/math","source_commit":"adc7f1241b42e322a6451854ab7e4b4c146bf78a","release_date":"2026-10-06","match":"full","resolution":"disproved","independently_reviewed":false,"second_pass":"agreed","families":[{"family":"155","title":"A counterexample to periodic tiling in dimension three","subject":"Combinatorics","manuscripts":[{"title":"A translational tile with no fully periodic tiling in dimension three","pdf_path":"preprints/A-translational-tile-with-no-fully-periodic-tiling-in-dimension-three-September-23-2026/paper.pdf","readme_path":"preprints/A-translational-tile-with-no-fully-periodic-tiling-in-dimension-three-September-23-2026/README.md","date":"2026-09-23","theorem_ref":"preprints/A-translational-tile-with-no-fully-periodic-tiling-in-dimension-three-September-23-2026/build/paper.tex#thm:main"}],"lean":{"doc_path":"lean/docs/155.md","comparators":["lean/ComparatorChallenges/PeriodicTilingThree.json","lean/ComparatorChallenges/PeriodicTilingThree.lean"],"declarations":["OAI.PeriodicTilingThree.periodic_tiling_counterexample_and_minimality"],"covers_main_theorem":true},"reasoning_trace":null}]}
\documentclass[11pt]{article}
\usepackage[margin=25mm]{geometry}
\usepackage{fontspec}
\setmainfont{Latin Modern Roman}
\usepackage{amsmath,amssymb,mathtools}
\usepackage{unicode-math}
\setmathfont{Latin Modern Math}
\usepackage{enumitem}
\usepackage{xurl,hyperref,bookmark}
\hypersetup{colorlinks=true,urlcolor=blue}
\setcounter{secnumdepth}{0}
\setlength{\parindent}{0pt}
\setlength{\parskip}{5pt}
\setlength{\emergencystretch}{3em}
\begin{document}
\section*{Periodic translational tiling in dimension three}

\subsection{OpenAI's claim}
\textit{Three-dimensional nonperiodic translational tile}: A finite lattice tile in $\mathbb Z^3$ is claimed to tile while admitting no tiling invariant under a finite-index translation group. Its thickening by unit cubes has the same obstruction in Euclidean space, including tilings by real translation vectors.

\subsection{Outline of the argument}
\textit{Three-dimensional nonperiodic translational tile}: A Sudoku-type word obstruction is encoded as several dependence and activation tests in a cyclic group. Stacking combines the tests into one tile, and an explicit model supplies a tiling. Marked component shapes rigidify any hypothetical fully periodic Euclidean tiling onto a grid, from which it would descend to a forbidden periodic solution of the encoded tests.

\subsection{Scope relative to this problem}
\textit{Three-dimensional nonperiodic translational tile}: The manuscript’s result has two distinct conclusions: a finite nonempty Z\textasciicircum{}3 tile with no finite-index-invariant complement, and its bounded positive-volume cube thickening with no full-rank-lattice-invariant real translation set despite an almost-everywhere tiling. The first clause supplies an exact counterexample to this question. The real conclusion explicitly permits arbitrary real translation vectors.

The claimed counterexample lies within the question’s hypotheses and contradicts its universal conclusion.

\subsection{Formal verification}
Family 155 describes Lean coverage for nonexistence of fully periodic lattice and Euclidean tilings for the three-dimensional construction.

\textit{Three-dimensional nonperiodic translational tile}: The manuscript links \texttt{PeriodicTilingThree.lean} (declaration \texttt{OAI.PeriodicTilingThree.periodic\_tiling\_counterexample\_and\_minimality}).

The documented selection includes this manuscript’s main assertion; that assertion retains the restrictions stated above.

This site has not run the Lean development. These links describe OpenAI’s selected formal statements and should be read with their stated scope.

\subsection{Status caveat}
OpenAI reports that an internal model produced the manuscripts. Their arguments have not been independently verified as proofs by this site. Any portion outside the documented formal scope may contain mistakes, and the catalogue label reports the announced mathematical claim.

\subsection{Sources}
\begin{itemize}
\item OpenAI, \emph{A translational tile with no fully periodic tiling in dimension three}, 2026-09-23; citation key \texttt{OAI:A-translational-tile-with-no-fully-periodic-tiling-in-dimension-three-September-23-2026}. \href{https://github.com/openai/math/blob/adc7f1241b42e322a6451854ab7e4b4c146bf78a/preprints/A-translational-tile-with-no-fully-periodic-tiling-in-dimension-three-September-23-2026/paper.pdf}{PDF (GitHub)}; \href{https://github.com/openai/math/raw/adc7f1241b42e322a6451854ab7e4b4c146bf78a/preprints/A-translational-tile-with-no-fully-periodic-tiling-in-dimension-three-September-23-2026/paper.pdf}{PDF (direct)}; \href{https://github.com/openai/math/blob/adc7f1241b42e322a6451854ab7e4b4c146bf78a/preprints/A-translational-tile-with-no-fully-periodic-tiling-in-dimension-three-September-23-2026/README.md}{Manuscript page and citation}; \href{https://github.com/openai/math/blob/main/preprints/A-translational-tile-with-no-fully-periodic-tiling-in-dimension-three-September-23-2026/paper.pdf}{Latest version}.
\item Family 155: \href{https://github.com/openai/math/blob/adc7f1241b42e322a6451854ab7e4b4c146bf78a/lean/docs/155.md}{Lean scope}.
\item \href{https://github.com/openai/math/blob/adc7f1241b42e322a6451854ab7e4b4c146bf78a/lean/ComparatorChallenges/PeriodicTilingThree.json}{Comparator PeriodicTilingThree.json}.
\item \href{https://github.com/openai/math/blob/adc7f1241b42e322a6451854ab7e4b4c146bf78a/lean/ComparatorChallenges/PeriodicTilingThree.lean}{Comparator PeriodicTilingThree.lean}.
\item \href{https://github.com/openai/math/blob/adc7f1241b42e322a6451854ab7e4b4c146bf78a/CONTENTS.md}{OpenAI catalogue entry}.
\end{itemize}
\end{document}
